% Results fragment; insert in the existing manuscript. No bibliography keys are invented.
\paragraph{M. Activity-based reference subgroup discovery}
The binary search returned 347 valid subgroups: 29 with one condition, 192 with two conditions, 126 with three conditions. 347 of the 347 reciprocity coefficients were negative; coefficients ranged from -1.150 to -0.216, against a global reference of -0.2413. The highest-quality subgroup was Sent$>$1 and Received$>$1 and Sent$<$=14 (n = 10,004,049; return rate 55.61\%; beta = -0.449; 95\% CI [-0.454, -0.444]). Its AME was -9.45 percentage points. The quality score was sqrt(n) times the absolute coefficient difference from the reference. 334/347 subgroup intervals excluded the reference point, which is reported as a descriptive diagnostic rather than a formal test of the subgroup-reference difference. Subgroup definitions used activity variables also included as controls.

\paragraph{A1. Reciprocity operationalization}
The beam search was repeated with proportional reciprocity as the target, with the same data, model, local discretization and settings. It returned 409 valid subgroups (409 with a negative coefficient), against 347 for binary reciprocity; 147 subgroups were visited by both searches. The two top-ten lists share 0 of 20 distinct descriptions (Jaccard overlap 0.000). Because thresholds are recomputed in every refined subgroup, the two searches follow different paths, so identical descriptions are a strict criterion. By observations, 95\% of those selected by the proportional top ten are also selected by the binary top ten (Jaccard overlap of the two selections 0.442). Each subgroup in either top-ten list was also estimated under the other measure: of the 20 that could be estimated under both, the coefficient sign agreed in 20 and the direction of the subgroup-reference difference in 20. Among all 147 subgroups visited by both searches, the direction of the difference agreed in 136. The weighted variant was not applicable: an exhaustive audit of the supplied raw networks found no repeated targets in any outgoing user-month list, so repeated-reply intensity was unavailable. Coefficients under the two measures refer to different contrasts and should not be interpreted as directly comparable probability changes.

\paragraph{A2. Dependence between observations}
User clustering increased the reference standard error by 13.1\% and top-ten subgroup standard errors by 2.0\%-5.3\% (median 3.9\%). All ten subgroup confidence intervals remained below zero under both approaches. Coefficients and the quality ranking were unchanged. M already used user-clustered inference; A2 supplies its explicit comparison with ordinary model-based uncertainty. A subgroup confidence interval excluding the global reference coefficient is a descriptive diagnostic, not a formal test of the difference between two estimated coefficients. The reference includes the subgroup, so a formal test must account for their joint uncertainty. Intervals are pointwise and have no multiple-testing or post-selection adjustment.

\paragraph{A3. Probability-scale comparison}
The reference AME was -4.44 percentage points. AMEs for the fixed top ten ranged from -10.14 to -7.25 percentage points. 10 of the 10 subgroup-reference differences had the same direction on the coefficient and probability scales. The Spearman correlation between rankings by absolute coefficient difference and absolute AME difference was -0.697, which shows how far the ordering of effect magnitudes depends on the scale. The most negative AME was -10.14 percentage points for Sent$>$1 and Received$>$1 and Sent$<$=8, rather than for the top-ten subgroup with the most negative coefficient. AMEs were evaluated over each model's own observations, and no AME confidence intervals were estimated.

\paragraph{Scope and interpretation}
These completed analyses use months 1-10, activity selectors and the two contemporaneous activity controls, with month and subreddit fixed effects. They do not include a community-context selector search or lagged controls. Results are descriptive rather than causal; pointwise intervals have no post-selection adjustment.
